% !Mode:: "TeX:UTF-8"
% 威海硕士学位论文开题报告的正文，照威海校区的开题报告范例写。括号里是范例的
% 填写要求，写完删掉。参考文献是第 7 节，由模板排在最后，这里不写；范例在那一节
% 的要求是：20 篇以上，外文不少于三分之一，近五年的一般不少于三分之一。

\section{课题来源及研究的目的与意义}

\subsection{课题的来源}

（注意：开题报告的字数应在 5000 字以上）

本课题来源于□□□项目“□□□□□□”。

\subsection{课题研究的背景和意义}

（注意：不少于 500 字）

发展国防工业、微电子工业等尖端技术需要精密和超精密的仪器设备，精密仪器设备要求高速、……

\section{国内外在该方向的研究现状及分析}

\subsection{国外研究现状}

1828年，R.R.Willis\cite{zhao1998}发表了一篇关于小孔节流平板中压力分布的文章，这是有记载的研究气体润滑的最早文献。……

\subsection{国内研究现状}

……\cite{lin1992,xin1999}

\subsection{国内外文献综述的简析}

（综合评述：国内外研究取得的成果，存在的不足或有待深入研究的问题。注意：不少于 500 字。）

……

\section{主要研究内容}

（注意：不少于 1000 字，撰写宜使用将来时态，切忌将论文目录直接作为研究内容，要突出本人研究内容）

……

\section{研究方案及进度安排，预期达到的目标和取得的研究成果}

\subsection{研究方案}

（注意：不少于 500 字）

……

\subsection{已完成的研究工作}

（注意：详细撰写目前已进行的研究工作内容和完成情况）

……

\subsection{进度安排和预期目标}

……

\section{为完成课题已具备和所需的条件和经费}

……

\section{预计研究过程中可能遇到的困难和问题，以及解决的措施}

\subsection{可能遇到的困难和问题}

……

\subsection{解决措施}

……
